% ============================================================================
%  Neural Network Field Map: Analysis Report
%  Author: G. Scriven (LHCb, Nikhef)
%  Date:   March 2026
% ============================================================================
\documentclass[11pt,a4paper]{article}

% ---- packages --------------------------------------------------------------
\usepackage[utf8]{inputenc}
\usepackage[T1]{fontenc}
\usepackage{lmodern}
\usepackage[margin=2.4cm]{geometry}
\usepackage{amsmath,amssymb,bm}
\usepackage{graphicx}
\usepackage{booktabs}
\usepackage{xcolor}
\usepackage{hyperref}
\usepackage{caption}
\usepackage{subcaption}
\usepackage{float}
\usepackage{enumitem}
\usepackage{fancyhdr}

% ---- colours ---------------------------------------------------------------
\definecolor{mlpblue}{HTML}{1F77B4}
\definecolor{nngreen}{HTML}{2CA02C}
\definecolor{siluorange}{HTML}{FF7F0E}
\definecolor{avxred}{HTML}{D62728}

% ---- hyperref setup --------------------------------------------------------
\hypersetup{
  colorlinks=true,
  linkcolor=mlpblue,
  citecolor=nngreen,
  urlcolor=siluorange,
}

% ---- headers ---------------------------------------------------------------
\pagestyle{fancy}
\fancyhf{}
\fancyhead[L]{\small NN Field Map Report}
\fancyhead[R]{\small LHCb, Nikhef}
\fancyfoot[C]{\thepage}

% ---- custom commands -------------------------------------------------------
\newcommand{\vB}{\vec{B}}
\newcommand{\qop}{q/p}
\DeclareMathOperator{\SiLU}{SiLU}
\DeclareMathOperator{\ReLU}{ReLU}

% ============================================================================
\begin{document}

% ---- title -----------------------------------------------------------------
\begin{center}
  {\LARGE\bfseries Neural Network Field Map:\\[4pt]
   Training, Deployment \& Benchmark Analysis}\\[12pt]
  {\large G.~Scriven}\\[4pt]
  {\normalsize LHCb, Nikhef}\\[4pt]
  {\normalsize March 2026}
\end{center}
\vspace{6pt}
\hrule
\vspace{12pt}

% ---- abstract ---------------------------------------------------------------
\begin{abstract}
This report documents the development and integration of a neural network
surrogate for the LHCb magnetic field map into the \texttt{TrackRungeKuttaExtrapolator}.
A grid search over 30 architectures (ReLU/SiLU $\times$ 1--3 layers $\times$
32--512 neurons) identified optimal accuracy--speed trade-offs. The deployed
model --- a single-layer ReLU network with 32 neurons (227 parameters,
416~FLOPs, 0.89~KB) --- replaces trilinear interpolation of a 4.5~MB field grid.
Three C++ inference paths were implemented: portable scalar, AVX2 SIMD, and
batched AVX2. RDTSC sub-operation profiling quantifies the field evaluation
fraction within the full Runge--Kutta pipeline and establishes the Amdahl's Law
ceiling for achievable speedup.
\end{abstract}

\tableofcontents
\newpage

% ============================================================================
\section{Introduction}
\label{sec:intro}

Track extrapolation in LHCb uses the Cash--Karp adaptive Runge--Kutta method
(4th/5th order) to propagate charged particles through the dipole magnetic
field. Each accepted step requires 6 field evaluations, and a typical
propagation from $z = 200$~mm to $z = 9500$~mm takes $\sim$11 steps.

The standard field lookup uses \textbf{trilinear interpolation} on an
$81 \times 81 \times 146$ grid (957,906 points, $\sim$4.5~MB in memory).
Profiling shows that field evaluation consumes $\sim$76\% of the total
propagation time, dominated by cache-miss latency from the large grid.

A neural network surrogate whose weights fit entirely in L1 cache
($\leq$1~KB) can eliminate this cache-miss bottleneck. The key requirements
are:
\begin{enumerate}[nosep]
  \item Accuracy: MAE $\lesssim 1$~Gauss across the full field volume
  \item Speed: $\leq$1500 FLOPs per evaluation (competitive with trilinear)
  \item Deployability: AVX2-friendly activation (no transcendental functions)
\end{enumerate}

% ============================================================================
\section{Grid Search: Architecture Exploration}
\label{sec:grid_search}

\subsection{Setup}

We trained 30 neural network architectures spanning:
\begin{itemize}[nosep]
  \item \textbf{Activations:} ReLU, SiLU
  \item \textbf{Depths:} 1, 2, 3 hidden layers
  \item \textbf{Widths:} 32, 64, 128, 256, 512 neurons per layer
\end{itemize}

All models map $(x, y, z) \to (B_x, B_y, B_z)$ with standardised
inputs/outputs. Training used:
\begin{itemize}[nosep]
  \item Optimizer: Adam, $\eta = 10^{-3}$
  \item Loss: MSE on normalised outputs
  \item Batch size: 4096
  \item Epochs: up to 200, early stopping (patience = 20)
  \item Validation: 100,000 held-out grid points (seed = 42)
  \item Data: \texttt{twodip.rtf} field map (957,906 grid points)
\end{itemize}

\subsection{Results}

Table~\ref{tab:grid_search} shows the top 10 architectures ranked by mean
absolute error (MAE) in Gauss.

\begin{table}[htbp]
\centering
\caption{Grid search: top 10 architectures by field MAE. Weight memory
  assumes float32 storage.}
\label{tab:grid_search}
\begin{tabular}{llrrrrr}
\toprule
Activation & Architecture & Params & FLOPs & MAE (G) & P99 (G) & Weight (KB) \\
\midrule
\multicolumn{7}{c}{\emph{(Fill from notebook output after running cell 2)}} \\
\bottomrule
\end{tabular}
\end{table}

Key observations:
\begin{itemize}[nosep]
  \item Accuracy improves monotonically with width at fixed depth.
  \item Deeper networks (2--3 layers) provide diminishing returns for the
        field map regression task.
  \item ReLU and SiLU achieve comparable MAE; ReLU is preferred for SIMD
        deployment (no \texttt{exp()} calls).
  \item The Pareto frontier (accuracy vs FLOPs) is dominated by shallow,
        wide architectures.
\end{itemize}

% ============================================================================
\section{Deployed Model: \texttt{field\_nn\_relu\_1L\_32H}}
\label{sec:deployed}

The model selected for deployment balances accuracy and inference speed:

\begin{center}
\begin{tabular}{ll}
\toprule
Property & Value \\
\midrule
Architecture     & $3 \to 32 \to 3$ (1 hidden layer, ReLU) \\
Parameters       & 227 (float32) \\
FLOPs            & 416 per evaluation \\
Weight memory    & 0.89 KB (fits in L1 cache) \\
Training epochs  & 200 (early stop at best epoch) \\
MAE              & 1.02 Gauss \\
P99 error        & 3.8 Gauss \\
Max error        & $\sim$20 Gauss \\
\bottomrule
\end{tabular}
\end{center}

\subsection{Normalization Constants}

Input and output standardisation (computed from the training grid):

\begin{center}
\begin{tabular}{lrrr}
\toprule
 & $x$ & $y$ & $z$ \\
\midrule
Input mean  & 0.0 & 0.0 & 6750.0 \\
Input std   & 2338.1 & 2338.1 & 4214.6 \\
\bottomrule
\end{tabular}
\qquad
\begin{tabular}{lrrr}
\toprule
 & $B_x$ & $B_y$ & $B_z$ \\
\midrule
Output mean & $\approx 0$ & 0.123 & $\approx 0$ \\
Output std  & 22.15 & 71.61 & 22.15 \\
\bottomrule
\end{tabular}
\end{center}

% ============================================================================
\section{C++ Integration}
\label{sec:cpp}

\subsection{Files Modified/Created}

\begin{table}[H]
\centering
\caption{Source files modified or created for the NN field map integration.}
\label{tab:files}
\begin{tabular}{lp{8cm}}
\toprule
File & Purpose \\
\midrule
\texttt{FieldMapNNReLUWeights.h} & Hardcoded weights + scalar, AVX2, and
  batched AVX2 inference functions \\
\texttt{FieldMapNNWeights.h} & SiLU variant (pre-existing) \\
\texttt{ExtrapolatorSubTimers.h} & RDTSC cycle-counter struct \\
\texttt{TrackFieldExtrapolatorBase.h/.cpp} & \texttt{UseNNFieldMap} Gaudi
  property + dispatch \\
\texttt{TrackRungeKuttaExtrapolator.cpp} & RDTSC sub-timer instrumentation \\
\texttt{TrackExtrapolatorTesterSOA.cpp} & ROOT ntuple sub-timer columns \\
\texttt{benchmark\_nn\_field\_map.py} & Gaudi benchmark config (4 variants) \\
\bottomrule
\end{tabular}
\end{table}

\subsection{Inference Functions}

Three inference paths are provided in \texttt{FieldMapNNReLUWeights.h}:

\begin{enumerate}
  \item \textbf{\texttt{evaluate\_relu(x,y,z)}} --- Portable scalar C++.
        Normalises input, computes $h = \ReLU(W_1 \cdot x_\mathrm{norm} + b_1)$,
        then $y = W_2 \cdot h + b_2$, and de-normalises output.

  \item \textbf{\texttt{evaluate\_relu\_avx2(x,y,z)}} --- AVX2 SIMD using
        \texttt{\_mm256\_fmadd\_ps} with pre-transposed weights. Processes
        8 neurons per vector operation.

  \item \textbf{\texttt{evaluate\_relu\_avx2\_batch<N>(...)}} --- Template
        function for track-level batching. Evaluates $N$ field points
        simultaneously, amortising weight-load overhead.
\end{enumerate}

\subsection{Runtime Configuration}

The field evaluation method is selected at runtime via the Gaudi property:
\begin{verbatim}
  ExtrapolatorInstance.UseNNFieldMap = "avx2_relu"
\end{verbatim}
Supported values: \texttt{""} (trilinear, default), \texttt{"scalar\_silu"},
\texttt{"scalar\_relu"}, \texttt{"avx2\_relu"}.

\subsection{Sub-Operation Profiling}

RDTSC cycle counters (\texttt{\_\_rdtsc()}) are inserted around five
sub-operations within each Runge--Kutta propagation:

\begin{enumerate}[nosep]
  \item \textbf{Field evaluation} --- magnetic field lookup or NN inference
  \item \textbf{Derivatives} --- Lorentz force computation
  \item \textbf{Butcher accumulation} --- Cash--Karp tableau weighted sums
  \item \textbf{Jacobian transport} --- covariance matrix propagation
  \item \textbf{Step-size control} --- error estimation and adaptive stepping
\end{enumerate}

These are recorded per-track in the \texttt{ExtrapolatorSubTimers} struct
and written to ROOT ntuples by \texttt{TrackExtrapolatorTesterSOA}.

% ============================================================================
\section{Production Benchmark}
\label{sec:benchmark}

\subsection{Configuration}

The benchmark (\texttt{benchmark\_nn\_field\_map.py}) runs four extrapolator
instances on identical tracks, propagating from $z = 3000$~mm to
$z = 12000$~mm through the dipole field:

\begin{enumerate}[nosep]
  \item \textbf{Trilinear} --- standard grid interpolation (baseline)
  \item \textbf{NN\_SiLU} --- scalar SiLU [32] inference
  \item \textbf{NN\_ReLU} --- scalar ReLU [32] inference
  \item \textbf{NN\_AVX2\_ReLU} --- AVX2 ReLU [32] inference
\end{enumerate}

Sub-operation timers are enabled (\texttt{EnableSubTimers = True}) to
capture per-track cycle breakdowns.

\subsection{Expected Output}

The benchmark produces a ROOT file with one ntuple per variant, containing:
\begin{itemize}[nosep]
  \item Per-track timing (\texttt{time} in ns)
  \item Final state residuals vs trilinear (\texttt{dx, dy, dtx, dty})
  \item Sub-operation cycle counts (\texttt{field\_cycles, deriv\_cycles, \ldots})
  \item Track properties (\texttt{qop, nsteps, nrejected})
\end{itemize}

\subsection{Analysis Plots}

The companion notebook (\texttt{nn\_field\_map\_analysis.ipynb}) generates:
\begin{itemize}[nosep]
  \item Wall-clock timing distributions and box plots
  \item Sub-operation horizontal bar charts per variant
  \item Stacked timing bar chart across variants
  \item Amdahl's Law: field evaluation fraction histogram
  \item Position/slope residual distributions vs trilinear
  \item Timing vs momentum ($|q/p|$) scatter plots
  \item LaTeX-ready summary tables
\end{itemize}

% ============================================================================
\section{Amdahl's Law Analysis}
\label{sec:amdahl}

If field evaluation accounts for a fraction $f$ of the total propagation
time, and the NN achieves a speedup $S_\mathrm{field}$ over trilinear for
field evaluation alone, then the overall speedup is bounded by:

$$
S_\mathrm{total} = \frac{1}{(1 - f) + f / S_\mathrm{field}}
$$

For the measured $f \approx 0.76$:
\begin{itemize}[nosep]
  \item $S_\mathrm{field} = 2\times$ $\Rightarrow$ $S_\mathrm{total} \leq 1.6\times$
  \item $S_\mathrm{field} = 5\times$ $\Rightarrow$ $S_\mathrm{total} \leq 2.6\times$
  \item $S_\mathrm{field} = \infty$ $\Rightarrow$ $S_\mathrm{total} \leq 4.2\times$
\end{itemize}

The sub-operation profiling from Section~\ref{sec:benchmark} provides the
actual $f$ for each variant, enabling precise Amdahl's Law predictions.

% ============================================================================
\section{Summary \& Outlook}
\label{sec:summary}

\subsection{Achievements}

\begin{enumerate}
  \item \textbf{Grid search} of 30 NN architectures identified the
        accuracy--speed Pareto frontier for field map surrogates.
  \item \textbf{Deployed model} (\texttt{field\_nn\_relu\_1L\_32H}):
        227 parameters, 416 FLOPs, MAE = 1.02~Gauss.
  \item \textbf{Three C++ inference paths}: scalar, AVX2, batched AVX2.
  \item \textbf{RDTSC sub-operation profiling} integrated into the full
        Runge--Kutta pipeline.
  \item \textbf{Benchmark infrastructure}: Gaudi config + ROOT ntuple output
        for systematic variant comparison.
\end{enumerate}

\subsection{Next Steps}

\begin{itemize}[nosep]
  \item Run production benchmark and fill in result tables.
  \item Investigate log-space weight training for improved accuracy.
  \item Profile on target hardware (Intel Xeon with AVX-512 support).
  \item Quantify the impact on track reconstruction efficiency and
        physics performance.
  \item Explore batch-level vectorisation across multiple tracks.
\end{itemize}

% ============================================================================
\end{document}
